% Einheit 3 von 4 – Lotfußpunkt (bank/abstaende/e3.jsonl, zusammenbau v0.1)
%% TODO Zeitmarke Einheit 3: Marken-Zeile nicht lesbar
%% TODO Zweigzeile Teil 1: was hier gelernt wird (Satz in Schülersprache aus dem Einheitstitel) – Einheit 3
\einheitenkopf[e3]{Einheit 3 von 4 · Lotfußpunkt}
\zweigzeile{Abitur GK}
\setcounter{aufgabe}{13}

%% TODO Titel: Ich-kann-Satz fehlt in der Bank (Platzhalter: Kettenname) – Nr. 14
%% TODO Anweisung: ein Satz über der ersten Teilaufgabe fehlt in der Bank – Nr. 14
\begin{aufgabe}{Lotfußpunkt}
%% TODO \rechenplatz{n}: Schrittzahl des Grundfalls fehlt in der Bank (nur bei fester Schreibform, 2.3 b)
\begin{teile}
\teil Was sagt die Gleichung $\overrightarrow{PF} \circ \vec r = 0$ in einem Lösungsweg? Kreuze an. \\ \kreuz{Punkt auf der Geraden} \\ \kreuz{Lotbedingung (Skalarprodukt null)} \\ \kreuz{Abstandsgleichheit}
\teil Lotfußpunkt F von P(4 | 2 | 3) auf g: $\vec x = (2 | 1 | 0) + t \cdot (2 | -1 | 2)$ – und der Abstand von P zu g? F( \leerfeld | \leerfeld | \leerfeld ), d ≈ \leerfeld
\teil An einem Zeltgestänge liefern die Gleichungen I $\overrightarrow{OQ} = \overrightarrow{OA} + r \cdot \overrightarrow{AB}$ und II $\overrightarrow{EQ} \circ \overrightarrow{AB} = 0$ gemeinsam einen Wert von r. Welche Bedeutung hat der Punkt Q für diesen Wert? \hfill (Abitur 2026 GK)
\teil Bei einem Carport soll von der Mitte M des Balkens AB eine möglichst kurze Strebe zum Balken CD führen; sie trifft CD im Punkt R. Beschreibe, wie man die Koordinaten von R ermitteln kann. \hfill (Abitur 2022 GK)
\teil Ein Holzklotz ist eine Pyramide über dem Quadrat ABCD mit A(2 | 1 | 0), B(6 | 1 | 0), C(6 | 5 | 0), D(2 | 5 | 0) und der Spitze E(2 | 5 | 3) senkrecht über D; 1 LE = 1 cm. B, D und E liegen in einer Symmetrieebene. Eine Ameise krabbelt auf kürzestem Weg von A über die Kante BE nach C. Wie lang ist der Weg? \leerfeld[cm]
\teil Die Gerade h: $\vec x = (3 | 3 | 3) + s \cdot (2 | -1 | 0)$ verläuft parallel zur Ebene E: $x + 2y + 2z = 6$. Gib eine Gleichung der Geraden in E an, die parallel zu h ist und von h den kleinsten Abstand hat. \hfill (Abitur 2022 LK) \\ $\vec x = $ \leerfeld
\teil Ein Mähroboter mit kreisförmiger Unterseite (Radius 0,3 m) fährt mit seinem Mittelpunkt auf der Geraden g mit Richtungsvektor (3 | 4 | 0) auf den Rand k der Rasenfläche zu; g trifft k im Punkt Q(6 | 8 | 0) unter dem Winkel 30°; 1 LE = 1 m. Der Roboter wendet, sobald seine Unterseite k berührt. Wo ist dann sein Mittelpunkt S? Nutze eine Skizze. \hfill (Abitur 2021 GK) \\ S( \leerfeld | \leerfeld | \leerfeld )
\teil Ein Dreieck ABC liegt in der xy-Ebene. Vorgelegt sind die Rechnungen $(4 | 2 | 0) \circ ((1 | 0 | 0) + \mu \cdot (4 | 2 | 0) - (2 | 3 | 0)) = 0 \Leftrightarrow \mu = 0{,}5$, S(3 | 1 | 0), und $\overrightarrow{OT} = \overrightarrow{OS} + |\overrightarrow{CS}| \cdot (0 | 0 | 1)$. Formuliere eine passende Aufgabenstellung und gib die Bedeutung von S an. \hfill (Abitur 2023 LK)
\end{teile}
\end{aufgabe}

%% TODO Titel: Ich-kann-Satz fehlt in der Bank (Platzhalter: Kettenname) – Nr. 15
%% TODO Anweisung: ein Satz über der ersten Teilaufgabe fehlt in der Bank – Nr. 15
\begin{aufgabe}{Lotfußpunkt – Fehler finden, Begründen, Anwendung}
\begin{teile}
\teil Finde den Fehler. Von M(3 | 3 | 4) soll eine möglichst kurze Stütze zum Balken SF mit S(0 | 1 | 6) und F(4 | 1 | 2) führen. Lukas nimmt den Mittelpunkt von SF als Auftreffpunkt: \rechnung{R &= (2 | 1 | 4)}
\teil Warum beschreibt das Gleichungspaar aus Geradengleichung und $\overrightarrow{PF} \circ \vec r = 0$ immer genau einen Lotfußpunkt? Begründe.
\teil Ein Karussellsitz ist im Punkt A(5 | 7 | 1,2) und dreht sich um die senkrechte Achse durch (2 | 3 | 0); 1 LE = 1 m. Welchen Weg legt der Sitz bei einer Umdrehung zurück? \leerfeld[m]
\end{teile}
\end{aufgabe}
